% Page 26 — Alban – XIVe siècle

\begin{center}
  {\Large\itshape Alban – XIV\textsuperscript{ème} siècle}
\end{center}
\vspace{1.2em}

\noindent
\lettrine{D}{e} Fonginières, les Anglais purent descendre jusqu’au « rec de la Capelanyé », ruisseau qui ne portait pas encore ce nom et qui coule en aval de la ferme de Capel le bas, sous la nouvelle gendarmerie, et sous le couvert des pentes boisées atteindre le village de Nougaret et donc le fameux rocher qui porterait comme unique trace de leur passage le nom de « Roc des Anglais ».

\medskip

Peut-être (ce mot est fort raisonnable) le calvaire de pierre qui se trouve à l’église et qui est classé monument historique, aurait-il été influencé dans sa décoration et son architecture par quelque descendant anglais ou plutôt irlandais demeuré dans le pays.

\medskip

\lettrine{E}{n} effet, cette croix a en son centre ce que les Anglais appellent la « rose des Tudors » et les divers tableaux qui sont sculptés dans le piédroit le sont à la façon des calvaires irlandais. Les pèlerins de Lourdes qui font le chemin de croix peuvent s’arrêter en fin de procession à un calvaire irlandais (de facture moderne).

\medskip

\lettrine{C}{ette} même « rose des Tudors » se trouve gravée sur une pierre en remploi – dans la façade de l’église de Miolles – Mais des motifs sculptés qui l’avoisinent portent la date de 1608.
